\section{Partie théorique} \label{sec:theorie}

\subsection{Experience 1, DMA} \label{subsec:osc_libres}

La DMA permet de caractériser le comportement viscoélastique d’un polymère soumis à une
contrainte sinusoïdale. L’instrument applique une contrainte périodique
\begin{equation}
    \sigma(t) = \sigma_0 \sin(\omega t),
\end{equation}
tandis que la déformation résultante est déphasée d’un angle $\delta$ :
\begin{equation}
    \varepsilon(t) = \varepsilon_0 \sin(\omega t + \delta).
\end{equation}

À partir de ce déphasage, les grandeurs viscoélastiques mesurées (exp1.pdf) sont :
\begin{align}
    E' &= \frac{\sigma_0}{\varepsilon_0}\cos\delta, \\
    E'' &= \frac{\sigma_0}{\varepsilon_0}\sin\delta, \\
    \tan\delta &= \frac{E''}{E'}.
\end{align}

Le pic du facteur d’amortissement $\tan\delta$ correspond à la transition vitreuse mécanique.
Comme cette transition dépend de la fréquence d’excitation, les températures des pics mesurés
sont exploitées à l’aide de la loi d’Arrhenius (exp1.pdf) :
\begin{equation}
    f = f_0 \exp\!\left(-\frac{E_a}{RT}\right),
\end{equation}
que l’on linéarise sous la forme
\begin{equation}
    \ln f = \ln f_0 - \frac{E_a}{R}\,\frac{1}{T}.
\end{equation}

Ainsi, la représentation de $\ln f$ en fonction de $1/T_g(f)$ permet d’extraire l’énergie
d’activation $E_a$ associée à la transition vitreuse du PMMA.


\subsection{Calorimétrie différentielle à balayage (DSC)}

La DSC mesure le flux de chaleur nécessaire pour maintenir un échantillon et une référence à
la même température lors d’une rampe thermique. Le signal correspond au flux enthalpique
(exp2.pdf) :
\begin{equation}
    \frac{dH}{dt}.
\end{equation}

À pression constante, la variation d’enthalpie est égale à la chaleur échangée (exp2.pdf) :
\begin{equation}
    dH = dQ_p,
\end{equation}
d'où
\begin{equation}
    \Delta H = Q_p.
\end{equation}

Le flux mesuré est relié à la capacité calorifique par
\begin{equation}
    \frac{dH}{dt} = m\,c_p\,\frac{dT}{dt}.
\end{equation}

Les transitions de premier ordre (fusion et cristallisation) apparaissent sous forme de pics
endothermiques ou exothermiques. La chaleur latente d’une transition se calcule par
\begin{equation}
    L = \frac{1}{m} \int_{\text{pic}} \frac{dH}{dt}\, dT.
\end{equation}

Pour un polymère semi-cristallin comme le PET, le taux de cristallinité s’obtient en comparant
la chaleur de fusion mesurée à celle d’un matériau parfaitement cristallin (exp2.pdf) :
\begin{equation}
    X_c = \frac{\Delta H_f - \Delta H_{c,\text{froide}}}{\Delta H_f^{100\%}}.
\end{equation}

Cette relation permet de déterminer quantitativement la proportion cristalline du PET à partir
des mesures DSC.

